\documentclass[11pt,a4paper]{article}
\usepackage[utf8]{inputenc}
\usepackage[T1]{fontenc}
\usepackage{lmodern}
\usepackage[margin=1in]{geometry}
\usepackage{amsmath,amssymb,amsthm}
\usepackage{booktabs}
\usepackage{longtable}
\usepackage{array}
\usepackage{microtype}
\usepackage[hidelinks]{hyperref}

\newtheorem{theorem}{Theorem}
\newtheorem{proposition}[theorem]{Proposition}
\newtheorem{lemma}[theorem]{Lemma}
\newtheorem{corollary}[theorem]{Corollary}
\theoremstyle{definition}
\newtheorem{definition}[theorem]{Definition}
\theoremstyle{remark}
\newtheorem{remark}[theorem]{Remark}

\newcommand{\lQ}{\lambda_{Q}}
\newcommand{\lR}{\lambda_{R}}
\newcommand{\lA}{\lambda_{A}}
\newcommand{\PsiL}{\Psi_{L}}
\newcommand{\src}[1]{{\footnotesize\texttt{#1}}}
\newcommand{\st}[1]{{\small\textsf{[#1]}}}

\title{\bfseries Certified Weil Positivity at $L = 4/5$:\\
Two-Sided Bounds, Parity Separation, and Cutoff Diagnostics}

\author{Elias De Jes\'us\\[2pt]
{\small Independent Researcher}\\[1pt]
{\small ORCID \href{https://orcid.org/0009-0007-0190-9143}{0009-0007-0190-9143}}\\[1pt]
{\small \texttt{dejesuselias10@gmail.com}}}

\date{November 2026 --- Technical Note (preprint), version 1.0.2}

\begin{document}
\maketitle

\begin{abstract}
\noindent
For the Weil explicit-formula quadratic form restricted to test functions
supported in a compact window $[-L,L]$ with $L = 4/5$, this note reports a
computer-assisted two-sided bound on its spectral infimum $\lQ(4/5)$ in one
fixed normalization:
\[
   1.41\times 10^{-17} \;\le\; \lQ(4/5) \;\le\; 1.9014\times 10^{-17}.
\]
The lower bound is obtained by a frequency-cutoff comparison at $T = 400$,
reduced to a $440$-dimensional Legendre block whose positive definiteness is
verified by a preconditioned congruence in Arb ball arithmetic, together with a
super-exponentially small complement estimate and a Schur-type coupling bound;
both parity sectors are certified. The upper bound is the Rayleigh quotient of
the \emph{original} form at one explicit even rational test function, evaluated
in the time domain through the autocorrelation and a compact archimedean
identity, so no oscillatory frequency tail is integrated. A separate certificate
gives $\lambda_{Q,\mathrm{odd}} \ge 2.3\times 10^{-17}$, which exceeds the upper
bound and therefore shows that the even infimum determines $\lQ$. Diagnostics at
cutoffs $T \in \{200,400,800\}$ measure the discarded exterior contribution and
bound the gain available from a larger cutoff. All constants are published as
midpoint--radius enclosures with a cold verifier. This is a single-window
computation: it establishes no arithmetic estimate, nothing about the location
of any zeta zero, and nothing about the Riemann hypothesis.
\end{abstract}

\section{Introduction}

\subsection{The object}

Weil's explicit formula pairs a test function against the zeros of $\zeta$, the
pole at $s=1$, the archimedean place and the primes. When the test function is
required to be supported in a fixed compact window, the resulting quadratic form
becomes a concrete object: a bounded-below Hermitian form on $L^2(-L,L)$ whose
positivity \emph{for every} $L$ is equivalent to the Riemann
hypothesis~\cite{Weil,Bombieri}, and whose positivity at one fixed $L$ is not. Finite-window Weil
positivity in this sense is the object studied in the spectral programs of
Connes and van Suijlekom~\cite{CvS}, Connes and Consani~\cite{CC21}, and
Connes, Consani and Moscovici~\cite{CCM}, and quantitatively in Zhu~\cite{Zhu}, which supplies the
finite-reduction strategy that the present computation follows.

The quantity computed here is the spectral infimum
\[
   \lQ(L) \;=\; \inf_{f} \frac{Q_{\zeta,L}(f)}{\|f\|_2^2}
\]
of that form at $L = 4/5$, in the normalization fixed in
Section~\ref{sec:setup}. Two things make a two-sided statement possible at this
window and difficult at larger ones. An \emph{upper} bound needs a single
admissible test function and is unconditional. A \emph{lower} bound needs the
whole window, hence a reduction of an infinite-dimensional positivity question
to a finite matrix plus certified tails; the cost of that reduction grows very
fast in $L$, which is the content of Theorem 1.4 of~\cite{Zhu}.

\subsection{What this note reports}

\begin{enumerate}
\item A certified lower bound $Q(f) \ge 1.41\times10^{-17}\|f\|_2^2$ at
  $L = 4/5$, in both parity sectors and hence for arbitrary complex tests
  (Section~\ref{sec:lower}).
\item A certified upper bound $\lQ(4/5) \le 1.9014\times10^{-17}$ from one
  explicit even rational test function, with the \emph{original} form evaluated
  in the time domain (Section~\ref{sec:upper}).
\item A certified parity separation: the odd sector satisfies
  $\lambda_{Q,\mathrm{odd}} \ge 2.3\times10^{-17}$, which is strictly above the
  upper bound, so the even infimum determines $\lQ$
  (Section~\ref{sec:parity}).
\item Cutoff diagnostics at $T \in \{200,400,800\}$ that measure what the
  comparison discards on fixed test functions, and a certified ceiling on what
  a larger cutoff can deliver (Section~\ref{sec:cutoff}).
\end{enumerate}

\subsection{Relation to prior work, and what is not claimed}
\label{sec:prior}

The finite-to-infinite reduction strategy --- a scalar gate on the exterior of a
frequency cutoff, a Legendre/Bessel representation of the interior, and
super-exponentially small tails --- is Theorem 1.1 of~\cite{Zhu}. Zhu describes
supplementary source code and certificates in~\cite[\S17]{Zhu}; these were not
obtained or executed in the present work, so his matrix step is neither re-run
nor verified here. The reduction was instead rederived in the normalization of
the present repository, and every constant below was computed from scratch. The bracket $8.9\times10^{-18} \le \lambda_{\min}(0.8) \le
2.27\times10^{-17}$ stated in the abstract of~\cite{Zhu} was used only as an
external comparison point and never as an input, an expected answer, or a
tolerance. The parity structure is likewise not new, and at this window it is
not merely the decoupling that is already published: \cite[Thm.~6.2]{Zhu}
certifies the \emph{separation} at support $1.6$ --- that is, at $L = 4/5$ --- and
does so with a far stronger odd-sector constant than the one obtained here,
concluding in addition that the ground state is simple and even. What
Section~\ref{sec:parity} contributes is an independent re-derivation of the
ordering from this pipeline's own certified data, nothing structural.

The computation continues the implementation developed alongside our
finite-window sensitivity study~\cite{Sensitivity}, whose repository supplies the
zeta-window normalization of Section~\ref{sec:setup}, the geometric-side
assembly, the pole conventions, and the closed-form archimedean evaluator used
as the validation reference in Section~\ref{sec:upper}(a). The even/odd pole
signs used here are the ones checked entrywise against that assembly in our
real-zero parity study~\cite{ParityNote}. Neither the $\delta^2$--$L^3$
sensitivity estimate of~\cite{Sensitivity} nor the planted-real-zero analysis
of~\cite{ParityNote} is an input to the lower bound proved here: the first
concerns a counterfactual perturbation of the zero set, and the second a
configuration that is nowhere assumed to exist. The finite-to-infinite reduction
follows~\cite{Zhu}; what is contributed here is its execution in certified
arithmetic, together with the resulting bracket, the numerical parity separation
of Section~\ref{sec:parity}, and the cutoff diagnostics of
Section~\ref{sec:cutoff}.

What follows is deliberately narrow, and the following are \emph{not} claimed.
No arithmetic estimate is obtained: nothing here improves any bound on primes,
on $\Lambda(n)$, on $\psi(x)$, or on a zero-free region. No statement is made
about the location of any zero of $\zeta$, and no range of zeros is verified.
No new positivity framework is introduced; the framework is that of the cited
sources, and the contribution is a certified computation inside it. The
computation is at one window, $L = 4/5$, and gives nothing at any larger support
--- there is no propagation theorem here, and a certificate at one $L$ does not
provide one. Nothing here bears on the Riemann hypothesis. This is a technical
note, not peer-reviewed research, and it is not claimed that anyone else uses or
has used this work.

\subsection{The same window, compared}
\label{sec:compared}

Because~\cite{Zhu} certifies at support $1.6$, its results and the present ones
describe the same object in the same normalization (its equations (2)--(3) are
\eqref{eq:Q}--\eqref{eq:Psi} below), so they can be put side by side. Doing so is
the fairest summary of what this note adds and what it does not.

\begin{center}
\begin{tabular}{lrr}
\toprule
quantity at $L = 4/5$ & \cite{Zhu}, certified & this note \\
\midrule
whole-window infimum
  & $[8.9\times10^{-18},\, 2.27\times10^{-17}]$
  & $[1.41\times10^{-17},\, 1.9014\times10^{-17}]$ \\
odd-sector lower bound
  & $8.206\times10^{-15}$
  & $2.3\times10^{-17}$ \\
second even min--max value
  & $\ge 2.085\times10^{-12}$
  & not certified here \\
simplicity of the ground state
  & \cite[Thm.~6.2]{Zhu}
  & not independently established here \\
\bottomrule
\end{tabular}
\end{center}

The one place this note improves on the published result is the width of the
whole-window bracket, which falls from $1.38\times10^{-17}$ to
$4.914\times10^{-18}$: approximately $64.4\%$ narrower, a reduction by a factor
of $2.81$. The present bracket is strictly nested inside the published one, both
endpoints moving inward. In the other three rows~\cite{Zhu} is stronger, and in
the odd sector by more than two orders of magnitude.

As a consistency check --- not as an endpoint, and not as a replay of his
certificate --- \cite{Zhu} also reports a numerical reference estimate
$\lambda_{\min}(Q|_{\text{window}}) = 1.656\times10^{-17}$ for this window, from
his own exact time-domain assembly. That value lies inside the bracket certified
here. It is a non-certified reference value in~\cite{Zhu} and is used here only
in that capacity.

\section{Setup, conventions, and the statement}
\label{sec:setup}

\subsection{Normalization}

Fix $L > 0$. For $f \in L^2(-L,L)$, extended by zero, write
\[
   F(t) \;=\; \int_{-L}^{L} f(u)\, e^{itu}\, du ,
   \qquad\text{so}\qquad
   \int_{\mathbb{R}} |F(t)|^2\, dt \;=\; 2\pi \|f\|_2^2 .
\]
The window Weil form used throughout is
\begin{equation}
\label{eq:Q}
   Q(f) \;=\; 2 s\, F(i/2)^2
     \;+\; \frac{1}{2\pi}\int_{\mathbb{R}} |F(t)|^2\, \PsiL(t)\, dt ,
\end{equation}
\begin{equation}
\label{eq:Psi}
   \PsiL(t) \;=\; H(t) - P_L(t), \qquad
   H(t) \;=\; \operatorname{Re}\psi\!\left(\tfrac14 + \tfrac{it}{2}\right) - \log\pi ,
   \qquad
   P_L(t) \;=\; \sum_{\log n < 2L} \frac{2\Lambda(n)}{\sqrt n}\,\cos(t \log n) .
\end{equation}
Here $\psi$ is the digamma function and $\Lambda$ is the von Mangoldt function.
The pole sign is $s = +1$ on the even sector and $s = -1$ on the odd sector; the
pole term is $2F(i/2)F(-i/2) = 2(C^2 - S^2)$ with $C = \int f \cosh(u/2)$ and
$S = \int f \sinh(u/2)$, so $C$ survives on even $f$ and $S$ on odd $f$, which
is exactly what the two signs encode. These signs are not asserted here: in the
parity study~\cite[\S3(c)]{ParityNote} the $\delta = 1/2$ real-pair perturbation
is compared entrywise with this repository's own pole term
$2\,\mathrm{pole\_sign}(\varepsilon)P_jP_k$ and agrees in all nine tested cells to
$6.22\times10^{-61}$, with $\mathrm{pole\_sign} = -1$ on the odd sector, while
the opposite sign convention is wrong by $O(1)$ to $O(10)$;
and~\cite[Remark 6]{ParityNote} identifies that $\delta = 1/2$ case as the parity
splitting of~\cite[Lemma 6.1]{Zhu}. The prime sum in~\eqref{eq:Psi} is
\emph{finite and exact}: the autocorrelation of a function supported in
$[-L,L]$ is supported in $[-2L,2L]$, so only $n$ with $\log n < 2L$ contribute
and there is no truncation. At $L = 4/5$ the active prime powers are
$\{2,3,4\}$.

This is the normalization of \src{src/zeta\_window.py} in the repository
accompanying this note, documented in~\cite{Sensitivity}, where the same object
is written $T_R = A - 2M$ with $A = \mathrm{Pole} + \mathrm{Arch} - (\log\pi)I$
and $2M = \mathrm{Prime}$. It is an imported convention, not derived here, and
the grouping of $-\log\pi$ is part of it. It was cross-checked against the repository's own
independent closed form for the constant test function
(Section~\ref{sec:upper}, item~(a)).

\subsection{The admissible class}

Since $H$ is even, continuous, and increasing on $t > 0$
(Lemma~\ref{lem:gate}), and $|P_L| \le A_L := \sum_{\log n < 2L}
2\Lambda(n)/\sqrt n < \infty$, the function $\PsiL$ is bounded below on
$\mathbb{R}$ and bounded on compacts. Hence for \emph{every} $f \in L^2(-L,L)$
the negative part of the integral in~\eqref{eq:Q} is finite and $Q(f)$ is
well defined in $(-\infty,+\infty]$. We therefore take the admissible class to
be all of $L^2(-L,L)$, with the convention that a lower bound
$Q(f) \ge \mu\|f\|_2^2$ is satisfied trivially where $Q(f) = +\infty$. This is a
choice of convention made in this note; it strengthens rather than weakens the
statement, and it removes any need for a density or form-domain argument.

\subsection{Three spectral quantities, kept apart}

Confusing the following three numbers is the easiest way to overstate a result
of this kind, so they are named separately and used separately throughout.

\begin{definition}
\label{def:three}
With $T > 0$ a frequency cutoff, $\beta$ a constant with $\PsiL(t) \ge \beta$
for $|t| \ge T$, and $V_N$ a finite-dimensional subspace of $L^2(-L,L)$:
\begin{align*}
  \lQ      &= \inf\{\,Q(f)/\|f\|_2^2 \,\}
             && \text{over the admissible class,}\\
  \lR(T)   &= \inf\{\,R_{T,\beta}(f)/\|f\|_2^2 \,\}
             && \text{over the same class, } R_{T,\beta} \text{ as in \eqref{eq:R},}\\
  \lA(T,N,\varepsilon) &= \lambda_{\min}\big(\text{compression of } R_{T,\beta} \text{ to } V_N\big)
             && \text{in the parity sector } \varepsilon .
\end{align*}
\end{definition}

\noindent Because $Q \ge R_{T,\beta}$ pointwise, and because an infimum over a
subspace is at least the infimum over the whole space,
\[
   \lR(T) \;\le\; \lQ , \qquad \lR(T) \;\le\; \lA(T,N,\varepsilon) ,
\]
and there is \textbf{no} inequality in either direction between $\lA$ and $\lQ$.
In particular a numerical least eigenvalue of a finite block is not an upper
bound for $\lQ$, and only the evaluation of $Q$ itself on an explicit test
function is.

\subsection{The theorem}

\begin{theorem}
\label{thm:main}
Let $L = 4/5$ and let $Q$ be the form~\eqref{eq:Q}--\eqref{eq:Psi} on
$L^2(-L,L)$.
\begin{enumerate}
\item[(i)] \emph{(Lower bound.)} For every $f \in L^2(-4/5,4/5)$,
  \[
     Q(f) \;\ge\; \tfrac{141}{10^{19}}\, \|f\|_2^2
     \;=\; 1.41\times 10^{-17}\, \|f\|_2^2 ,
  \]
  in both parity sectors separately and hence for arbitrary complex $f$.
\item[(ii)] \emph{(Upper bound.)} There is an explicit even test function
  $f_{400}$, a Legendre polynomial of degree $438$ with rational coefficients of
  denominator $10^{40}$, such that
  \[
     \lQ(4/5) \;\le\; \frac{Q(f_{400})}{\|f_{400}\|_2^2}
     \;\le\; U \;=\; 1.901382911145464406030000\times10^{-17}
     \;<\; 1.9014\times10^{-17}.
  \]
\item[(iii)] \emph{(Parity separation.)} For every odd
  $f \in L^2(-4/5,4/5)$, $Q(f) \ge 2.3\times10^{-17}\|f\|_2^2$. Since
  $f_{400}$ is even and $2.3\times10^{-17} - U > 3.98\times10^{-18} > 0$,
  \[
     \inf_{f \text{ even}} \frac{Q(f)}{\|f\|_2^2} \;\le\; U
     \;<\; 2.3\times10^{-17} \;\le\;
     \inf_{f \text{ odd}} \frac{Q(f)}{\|f\|_2^2} ,
  \]
  so the even infimum equals $\lQ(4/5)$.
\end{enumerate}
\end{theorem}

Parts (i) and (iii) are computer-assisted: they rest on interval arithmetic in
the Arb midpoint--radius model~\cite{Arb}, including all analytic error terms, on exact rational input data, and
on a cold verifier that recomputes every decisive implication. Part (ii) needs
one test function and is unconditional given the normalization.

\begin{remark}
Theorem~\ref{thm:main}(iii) identifies which sector carries the infimum. It does
\emph{not} assert that the infimum is attained, that a minimizer is unique, or
that the least eigenvalue is simple. None of those is established here.
\end{remark}

\begin{remark}
The bracket width is $1.9014\times10^{-17} - 1.41\times10^{-17} < 4.92\times
10^{-18}$, about $35\%$ of the upper endpoint. Section~\ref{sec:cutoff} measures
where that width comes from and bounds how much of it the present method can
remove.
\end{remark}

\section{The certified lower bound}
\label{sec:lower}

The lower bound is assembled in five steps. Steps 1--5 are proved here in the
normalization of Section~\ref{sec:setup}; the strategy is that
of~\cite[Thm.~1.1]{Zhu}.

\subsection{Step 1: the exterior gate}

\begin{lemma}
\label{lem:gate}
$t \mapsto \operatorname{Re}\psi(\tfrac14 + \tfrac{it}{2})$ is strictly
increasing on $t > 0$.
\end{lemma}

\begin{proof}
From the digamma series,
\[
   \operatorname{Re}\psi\!\left(\tfrac14+\tfrac{it}{2}\right)
   = -\gamma + \sum_{n \ge 0}
     \left[\frac{1}{n+1} - \frac{n + 1/4}{(n+1/4)^2 + t^2/4}\right],
\]
and each summand is strictly increasing in $t > 0$. Termwise differentiation is
justified by local uniform convergence.
\end{proof}

Consequently, with $A_L = \sum_{\log n < 2L} 2\Lambda(n)/\sqrt n$ and
$\beta := H(T) - A_L$, we get $\PsiL(t) \ge \beta$ for all $|t| \ge T$, because
$P_L \le A_L$ pointwise. At $L = 4/5$,
\[
   A_L = 2.941973525223620455503909 \pm 3.91\times10^{-25},
\]
and the gate values at the three cutoffs used in this note are
\begin{center}
\begin{tabular}{lll}
\toprule
$T$ & $\beta = H(T) - A_L$ & status \\
\midrule
$200$ & $0.5184657332438467199873536 \pm 2.62\times10^{-26}$ & certified \\
$400$ & $1.211613695058064549471502 \pm 1.89\times10^{-26}$ & certified \\
$800$ & $1.904761070930776888621206 \pm 1.68\times10^{-25}$ & certified \\
\bottomrule
\end{tabular}
\end{center}
with the increments $\beta_{400} - \beta_{200} = 0.6931479618142178294841484 \pm
7.24\times10^{-27}$, numerically close to $\log 2$ as the asymptotics of $H$
suggest. All three are \emph{positive}, so the set
$\{|t| \ge T : \PsiL(t) < \beta\}$ is empty and no occupation or measure
estimate on a bad exterior set is needed anywhere below.

\subsection{Step 2: the comparison form}

Because $\int_{|t| \ge T} (\PsiL - \beta)|F|^2 \ge 0$, equation~\eqref{eq:Q}
gives $Q(f) \ge R_{T,\beta}(f)$ with
\begin{equation}
\label{eq:R}
   R_{T,\beta}(f) \;=\; \beta \|f\|_2^2 \;+\; 2 s F(i/2)^2
     \;+\; \frac{1}{2\pi}\int_{-T}^{T} \big(\PsiL(t) - \beta\big) |F(t)|^2\, dt .
\end{equation}
The discarded quantity is exactly
\begin{equation}
\label{eq:e}
   e_T(f) \;:=\; \frac{Q(f) - R_{T,\beta}(f)}{\|f\|_2^2}
   \;=\; \frac{1}{2\pi\|f\|_2^2}\int_{|t| > T} \big(\PsiL(t) - \beta\big)|F(t)|^2\, dt
   \;\ge\; 0 ,
\end{equation}
which is measured in Section~\ref{sec:cutoff}.

\begin{proposition}[monotonicity in $\beta$]
\label{prop:beta}
By Plancherel, $R_{T,\beta}(f) = 2sF(i/2)^2 + \frac{1}{2\pi}\int_{|t|\le T}
\PsiL |F|^2 + \frac{\beta}{2\pi}\int_{|t| > T}|F|^2$, so for
$\beta_2 \ge \beta_1$,
\[
   R_{T,\beta_2}(f) - R_{T,\beta_1}(f)
   = \frac{\beta_2-\beta_1}{2\pi}\int_{|t|>T}|F|^2 \;\ge\; 0 .
\]
On the finite compression this reads $A(\beta_2) - A(\beta_1) =
(\beta_2-\beta_1)(I - K_T)$, where $K_T$ is the compression of the Fourier
energy-concentration operator and $0 \le K_T \le I$.
\end{proposition}

Proposition~\ref{prop:beta} is recorded because it rules out a tempting move:
\emph{lowering} $\beta$ cannot raise the exact comparison form, its infimum, or
the least eigenvalue of its compression to a fixed subspace. Inspecting only the
factor $\PsiL - \beta$ inside the integral misses the $\beta\|f\|_2^2$ term. The
identity was checked numerically with maximum residual $5.23\times10^{-143}$ and
diagonal entries of $K_T$ in $[0.9420888637695,\, 0.9980132600038]$ at $T=200$,
$N=16$.

\subsection{Step 3: the Legendre basis}

Let $\varphi_n(u) = \sqrt{(2n+1)/(2L)}\, P_n(u/L)$ be the orthonormal Legendre
basis of $L^2(-L,L)$. Then
\[
   F_n(t) \;=\; i^n \sqrt{2L(2n+1)}\; j_n(Lt) ,
   \qquad \sum_{n \ge 0}(2n+1) j_n(x)^2 = 1 ,
\]
with $j_n$ the spherical Bessel function. The form is assembled in the real
coordinates $y_n = \sigma_n x_n$ with $\sigma_n = (-1)^{\lfloor n/2 \rfloor}$,
which makes $v_n(t) = \sqrt{2L(2n+1)}\,j_n(Lt)$ the real working transform. This
basis is used, rather than the cosine basis of the repository's original
assembly, because $j_n(x)$ is super-exponentially small for $n > ex/2$, which is
what makes Step 4 effective; a cosine basis gives only $O(n^{-2})$ tail decay
and cannot reach the required magnitude at an acceptable dimension.

\subsection{Step 4: the complement}

For $f \in V_N^{\perp}$ with $\|f\|_2 = 1$, Cauchy--Schwarz in $n$ gives
$|F(t)|^2 \le 2L\, S_N(tL)$ with $S_N(x) = \sum_{n \ge N}(2n+1)j_n(x)^2$. Since
$S_N$ is increasing on $0 \le x < N$,
\[
   g_N \;:=\; \sup_{f \in V_N^{\perp},\, \|f\|=1}
     \frac{1}{2\pi}\int_{-T}^{T}|F|^2
   \;\le\; \frac{2LT}{\pi}\, S_N(TL) .
\]
This is a lossy but rigorous surrogate for the $(N+1)$-st Slepian
eigenvalue~\cite{Slepian},
and it is super-exponentially small past the Shannon number $2TL/\pi$. At
$T = 400$, $N = 440$:
\[
   S_N(320) = 9.260051461963608\times10^{-63},
   \qquad
   g_N = 1.886442193224756\times10^{-60}.
\]

\subsection{Step 5: the Schur-type block bound, with strict tests}

Split $f = f_N + f_{\perp}$. With $a$ a certified lower bound for the least
eigenvalue of the compression of $R_{T,\beta}$ to $V_N$, $d$ a certified lower
bound for $R_{T,\beta}$ on $V_N^{\perp}$, and $\kappa$ a certified bound on the
coupling, a $2\times2$ comparison yields positivity as soon as $ad > \kappa^2$.
The threshold $\kappa^2/d$ is a \emph{break-even point}, not the resulting
bound. To avoid the cancellation-prone explicit eigenvalue formula, the final
implication is tested as three \emph{strict} interval signs for the proposed
constant $\mu$:
\begin{equation}
\label{eq:schur}
   a > \mu , \qquad d > \mu , \qquad (a - \mu)(d - \mu) > \kappa^2 .
\end{equation}
Equality anywhere in~\eqref{eq:schur} would give only a nonnegative comparison,
so strictness is required and is checked as a sign of an interval, never of a
midpoint. The coupling constant is
$\kappa = 2|p_N||p_{\perp}| + M\sqrt{g_N}$ with
$M = \sup_{|t| \le T}|\PsiL - \beta|$ and $p$ the pole vector
$p_n = \sqrt{2L(2n+1)}\, i_n(L/2)$ taken in the same $y$-coordinates; $|p|^2 =
\int_{-L}^{L} e^{-u}\,du = 2\sinh L$ exactly. $M$ is obtained from the
monotonicity of $H$ alone, with no covering argument.

\subsection{Verified positive definiteness}

A naive interval $LDL^{\mathsf T}$ factorization is unusable at this dimension:
entry radii of order $10^{-38}$ amplify by roughly $6.5$ per pivot and reach
$O(1)$ by step $47$. Instead, positive definiteness of $B = A - aI$ is
established by a preconditioned congruence in the sense of Rump~\cite{Rump}. Let $R$ be the
upper triangular Cholesky factor of the \emph{midpoint} of $B$, treated as exact
input, and set $C = R^{-1}$ and
\[
   G \;=\; C^{*} B\, C \;=\; R^{-\mathsf T}\,(A - aI)\, R^{-1},
\]
computed by two explicit forward substitutions in ball arithmetic. The verified
inequality is
\[
   \| I - G \|_{\infty} \;<\; 1 .
\]
This proves $B \succ 0$: the bound places every eigenvalue of the Hermitian $G$
in $(0,2)$, so $G \succ 0$; if $C$ were singular then $G = C^{*}BC$ would be
singular, so $C$ is nonsingular; and then for $x \ne 0$, writing $x = Cy$ with
$y \ne 0$, $x^{*}Bx = y^{*}Gy > 0$. It is a congruence, not a one-sided inverse
residual --- the latter would establish nonsingularity only.

\subsection{Quadrature}

The interior integral in~\eqref{eq:R} is evaluated by panelled
Gauss--Legendre quadrature with certified nodes and weights (FLINT's
\src{legendre\_p\_root}). The symbol, prime comb and the full basis-transform
vector are evaluated once per node and reused across all matrix entries, so no
entry is integrated adaptively. The quadrature remainder is bounded by a
Bernstein-ellipse estimate on genuinely holomorphic extensions: the half-sum
\[
   H_{\mathrm{hol}}(z) \;=\; \tfrac12\left[
     \psi\!\left(\tfrac14 + \tfrac{iz}{2}\right)
   + \psi\!\left(\tfrac14 - \tfrac{iz}{2}\right)\right] - \log\pi
\]
agrees with $H$ on the real axis and has its nearest poles at $z = \pm i/2$
(note that $\operatorname{Re}\psi(\frac14 + \frac{iz}{2})$ is \emph{not} a
holomorphic extension), and the basis product is extended as $v_n(z)v_m(-z)$
rather than by conjugation at a complex argument. On the ellipses, $|v_n(z)|$ is
bounded uniformly in $n$ by
\begin{equation}
\label{eq:cs}
   \Big|\int_{-L}^{L}\varphi(u) e^{izu}\,du\Big|
   \;\le\; \Big(\int_{-L}^{L} e^{-2u \operatorname{Im} z}du\Big)^{1/2}
   \;\le\; \sqrt{2L}\, e^{L|\operatorname{Im} z|}
\end{equation}
for any $L^2$-normalized $\varphi$ supported in $[-L,L]$. Inequality
\eqref{eq:cs} is immediate from Cauchy--Schwarz and is \emph{not} claimed as
new; it is recorded because it replaces a direct interval evaluation of
$j_n$ on wide complex balls, which is numerically useless here. The remainder is
folded into the matrix entries as a widening exactly once.

\subsection{The $T = 400$ certificate}

\begin{center}
\begin{tabular}{ll}
\toprule
parameter & value \\
\midrule
window & $L = 4/5$ \\
cutoff & $T = 400$ \\
dimension & $N = 440$ (Legendre degrees $0..439$; $220$ modes per parity) \\
quadrature & $800$ panels of half-length $1/4$, Gauss--Legendre order $64$ \\
Bernstein data & $b_e = 0.30$ (inside the strip $|\operatorname{Im} z| < 1/2$), $\rho = 2.762049935$ \\
remainder factor & $3.551349410\times10^{-57}$ \\
precision & $900$ bits (Bessel), $384$ bits (accumulation), $1400$ bits (verification) \\
\bottomrule
\end{tabular}
\end{center}

$N = 440$ was chosen from the actual inequalities, not from the heuristic
$eLT/2 = 434.9$. The certified constants are
\begin{align*}
  \beta &= 1.211613695058064549471502, &
  M &= 9.525770639507350587208369, \\
  \kappa &= 1.308343771285388259029018\times10^{-29}, &
  \kappa^2/d &= 1.412796364751588\times10^{-58},
\end{align*}
with pole tail $|p_{\perp}| = 7.993372152149564\times10^{-1282}$.

\begin{center}
\begin{tabular}{lccccc}
\toprule
sector & $a$ (certified) & $\|I-G\|_{\infty}$ & $d$ & $\mu$ & margin in \eqref{eq:schur} \\
\midrule
even & $1.42\times10^{-17}$ & $1.4973903\times10^{-14}$ & $1.211613695058065$ & $1.41\times10^{-17}$ & $1.21161369506\times10^{-19}$ \\
odd  & $2.4\times10^{-17}$  & $4.3535838\times10^{-22}$ & $1.211613695058065$ & $2.3\times10^{-17}$  & $1.21161369506\times10^{-18}$ \\
\bottomrule
\end{tabular}
\end{center}

The whole-window constant is the smaller of the two sector constants, so
$\mu = 1.41\times10^{-17}$ and \emph{the even sector binds}. Assembly took
$3532$ s, the full run $3581$ s, with peak resident memory $393.4$ MiB; the two
Rump factorizations at $N = 440$ took under $50$ s together.

\subsection{Parity and complex tests}

The form decouples by parity: the interior matrix has $C_{nm} = 0$ when $n+m$ is
odd, because $v_nv_m$ is then odd on a symmetric interval, and the pole term
$2(C^2 - S^2)$ separates into its even and odd parts. For complex $f = u + iv$
the Hermitian form gives $Q(f) = Q(u) + Q(v)$, and each real part splits by
parity. Hence the two sector certificates cover the full class, which is what
makes Theorem~\ref{thm:main}(i) a statement about arbitrary complex $f$.

\begin{remark}
A positive-definiteness check cannot serve as an oracle for the pole sign. The
direct convention checks of Section~\ref{sec:upper} are retained for that
purpose, because an earlier version of this computation carried a pole-phase
error that produced a plausibly negative pivot
(Appendix~\ref{app:errors}, item~1).
\end{remark}

\section{The original-form upper bound}
\label{sec:upper}

\subsection{Why the time domain}

A variational upper bound requires $Q$ itself, not $R_{T,\beta}$. Evaluating
\eqref{eq:Q} directly would mean certifying an oscillatory integral over all of
$\mathbb{R}$. Instead $Q$ is evaluated on the geometric side via the
autocorrelation
\[
   g(x) \;=\; \int f(s+x) f(s)\, ds ,
\]
supported in $[-2L,2L]$, together with the compact archimedean identity
(\cite[Lemma 2.5]{Zhu}, transcribed into the present normalization)
\begin{equation}
\label{eq:arch}
   \mathrm{Arch}(g) = -\big[\gamma_E + \log\pi + \log(1 - e^{-4L})\big] g(0)
   + \int_{0}^{2L} \frac{2\big[e^{-2x}g(0) - e^{-x/2}g(x)\big]}{1 - e^{-2x}}\, dx ,
\end{equation}
\begin{equation}
\label{eq:Qcompact}
   Q(f) = 2sF(i/2)^2 + \mathrm{Arch}(g)
   - \sum_{\log n < 2L} \frac{2\Lambda(n)}{\sqrt n}\, g(\log n) .
\end{equation}
This is the route by which~\cite[\S6.2]{Zhu} evaluates its own variational upper
bounds at this window, there with shift-correlation integrands of degree at most
$239$ integrated exactly by a $160$-node Gauss rule. Two features make it cheap
and rigorous. They are set out below because the present error accounting
depends on them, not as new observations.

\paragraph{Exact inner quadrature.} If $f$ is a polynomial of degree $d$ on
$[-L,L]$ then $g$ is a polynomial of degree at most $2d$ on $[0,2L]$, and the
integrand $f(s+x)f(s)$ has degree at most $2d$ in the integration variable, so
an $m$-node Gauss--Legendre rule is exact once $2m - 1 \ge 2d$, i.e.\
$m \ge d+1$. The implementation uses $m = d+2$; for $d = 438$ that is $440$
nodes. \emph{Exactness} here means no polynomial quadrature truncation error in
exact arithmetic --- node, weight, coefficient and rounding enclosures are still
carried as ball radii, and the outer integral in~\eqref{eq:arch} keeps its own
certified Bernstein remainder.

\paragraph{The removable endpoint, by an identity.} The integrand
of~\eqref{eq:arch} is removable at $x = 0$ and its limit involves
$c = -g'(0^{+})$. For even $f$,
\[
   g'(x) \;=\; -f(L) f(L-x) + \int_{-L}^{L-x} f'(s+x) f(s)\, ds ,
\]
and at $x = 0$ the integral is $\tfrac12[f^2]_{-L}^{L} = 0$, so
\begin{equation}
\label{eq:c}
   c \;=\; -g'(0^{+}) \;=\; f(L)^2 \qquad\text{exactly,}
\end{equation}
and the integrand limit is $c - \tfrac32 g(0)$. This closed form is an
implementation convenience obtained by elementary differentiation and
integration; \cite[Lemma 2.5]{Zhu} leaves $c$ as the one-sided limit and remarks
only that it is nonzero in general. Its absence from the displayed formulas
there is not evidence that it is new, and no novelty is claimed for it. The outer integrand is analytic
in $|\operatorname{Im} x| < \pi$; panels are geometric near the origin so that no
Bernstein ellipse encloses it, and the residual head $[0,10^{-40}]$ is bounded
by an explicit constant.

\subsection{Validation of the representation}

\begin{enumerate}
\item[(a)] On the constant $f = \varphi_0$, the compact route gives
  $Q(\varphi_0) = 0.080274054959269611214 \pm 3.09\times10^{-22}$, matching the
  repository's independent closed form $Q[0,0] = 0.080274054959269611214$ to all
  $21$ printed digits (difference $-3.08\times10^{-22}$). That closed form is the
  evaluator of~\cite{Sensitivity}, which sums the archimedean series analytically
  through the digamma, Lerch $\Phi$ and Hurwitz $\zeta$ functions; it is an
  independent route to $Q$, not an independent implementation of the present
  one.
\item[(b)] A degree-$0$ check cannot see a pole-phase error, since
  $\sigma_0 = +1$. Checks at $\varphi_2$ and at the signed mixed function
  $\varphi_0 - \varphi_2$ were therefore added, with the pole moment computed
  \emph{directly} from the physical-space function rather than from the
  pole-vector helper, and with the finite prime shifts re-enumerated
  independently of the symbol helper:
  \begin{center}
  \begin{tabular}{lll}
  \toprule
  $f$ & $Q(f)$ & pole moment (direct) \\
  \midrule
  $\varphi_0$ & $0.080274054959269611214$ & $1.29891290374844725$ \\
  $\varphi_2$ & $0.389336213205013655883$ & $0.030516224225381024$ \\
  $\varphi_0 - \varphi_2$ & $0.117009437226589934780$ & $1.26839667952306623$ \\
  \bottomrule
  \end{tabular}
  \end{center}
  All agree with the helper values, the first to $10^{-25}$.
\item[(c)] \emph{Disclosure.} This is a second \emph{representation} of the same
  form, not an independent implementation: it shares the Arb library, the symbol
  helper for the comb data, and the pole-vector helper. It is not described as an
  independent check.
\end{enumerate}

\subsection{The frozen test functions and the bound}

Two trial functions are used, each a frozen exact rational vector in the
Legendre $y$-coordinates with denominator $10^{40}$:
\begin{itemize}
\item $f_{200}$: $110$ even modes, degrees $0..218$, $\|f\|_2^2 =
  2.275244443650965360970550$;
\item $f_{400}$: $220$ even modes, degrees $0..438$.
\end{itemize}
Both were \emph{located} by an exploratory floating-point eigensolver, then
rounded to exact rationals and certified; the eigensolver contributes no
certified claim, and in particular supplies no eigenvalue.

For $f_{200}$, equation~\eqref{eq:Qcompact} gives, at $94$ outer panels and head
cut $10^{-40}$,
\[
  \text{pole} = 3.2303213242823875400,\qquad
  \mathrm{Arch} = -3.1402989129419417722,
\]
\[
  \text{comb} = 0.090022411340445684793,
\]
whose signed sum is $8.30143811294573721\times10^{-17}$: a cancellation of about
sixteen digits, which is why the evaluation is done at high precision in ball
arithmetic and why the compact representation --- with its exact inner
quadrature and the exact identity~\eqref{eq:c} --- matters. Dividing by the norm,
\[
  q(f_{200}) = 3.64859175290407702\times10^{-17} \pm 4.15\times10^{-35}.
\]
This is above the then-current upper bound and so does not improve it; the
expected outcome, since $f_{200}$ minimizes a comparison form, not $Q$.

For $f_{400}$,
\[
  q(f_{400}) = 1.90138291114546440\times10^{-17} \pm 6.03\times10^{-35},
\]
whose certified upper endpoint is
\[
  U \;=\; 1.901382911145464406030000\times10^{-17},
\]
giving Theorem~\ref{thm:main}(ii).

\begin{remark}[on quoting this number]
\label{rem:endpoint}
The decimal $1.90138291114546440\times10^{-17}$ is the \emph{midpoint} of the
enclosure and is \emph{not} a valid strict upper bound for it; the strict upper
endpoint is $U$ above. Earlier internal drafts of this work quoted the midpoint
as the bracket's upper endpoint, and a separate internal record quoted a looser
endpoint $1.9013829111454646892\times10^{-17}$ arising from a low-precision
re-parse of the serialized ball; both are corrected
(Appendix~\ref{app:errors}, item~6). The rounded value $1.9014\times10^{-17}$
used in the abstract and in Theorem~\ref{thm:main} is above every one of them
and is the safe form in which to quote the result.
\end{remark}

\section{The parity separation}
\label{sec:parity}

The $T = 400$ certificate of Section~\ref{sec:lower} delivers $a = 2.4\times
10^{-17}$ in the odd sector against $a = 1.42\times10^{-17}$ in the even sector.
Re-evaluating the three strict signs~\eqref{eq:schur} from the exact certified
data --- not from rounded report values --- at $\mu_{\mathrm{odd}} =
2.3\times10^{-17}$ gives
\[
   a > \mu_{\mathrm{odd}},\qquad d > \mu_{\mathrm{odd}},\qquad
   (a - \mu_{\mathrm{odd}})(d - \mu_{\mathrm{odd}}) > \kappa^2 ,
\]
with margin $1.21161369506\times10^{-18}$. So the whole odd sector satisfies
$Q(f) \ge 2.3\times10^{-17}\|f\|_2^2$. Combining with $U$ from
Section~\ref{sec:upper}, and using that $f_{400}$ is even,
\[
   \mu_{\mathrm{odd}} - U \;=\; 3.986170888545355939700000\times10^{-18}
   \;>\; 0 ,
\]
so
\[
   \inf_{\text{even}} \frac{Q(f)}{\|f\|_2^2} \;\le\; U \;<\;
   \mu_{\mathrm{odd}} \;\le\; \inf_{\text{odd}} \frac{Q(f)}{\|f\|_2^2},
\]
which with the orthogonal parity decomposition identifies the even infimum as
$\lQ(4/5)$. The global lower bound remains $1.41\times10^{-17}$, governed by the
even sector. As noted in Section~\ref{sec:prior}, none of this structure is new: at this very
window \cite[Thm.~6.2]{Zhu} certifies the same ordering with
$\lambda^{\mathrm{odd}}_1 \ge 8.206\times10^{-15}$, more than two orders of
magnitude above the constant obtained here. What is recorded above is only that
this pipeline's own certified numbers separate the two sectors. Nothing about
simplicity, uniqueness or attainment follows from them --- simplicity and
evenness at this window were previously established by~\cite[Thm.~6.2]{Zhu}.

\section{Cutoff diagnostics}
\label{sec:cutoff}

\subsection{The cutoff comparison, proved}

\begin{proposition}
\label{prop:cutoff}
Let $\PsiL \ge \beta_i$ on $|t| > T_i$ for $i = 1,2$, with $T_2 \ge T_1$ and
$\beta_2 \ge \beta_1$. Then
\[
   R_{T_2,\beta_2}(f) - R_{T_1,\beta_1}(f)
   = \frac{1}{2\pi}\int_{T_1 < |t| \le T_2} (\PsiL - \beta_1)|F|^2
   + \frac{\beta_2 - \beta_1}{2\pi}\int_{|t| > T_2} |F|^2
   \;\ge\; 0 ,
\]
both terms being nonnegative. Hence $\lR(T_2) \ge \lR(T_1)$ for the exact forms.
\end{proposition}

Two cautions. This does not imply that a \emph{certified numerical} bound
improves, since error terms change with the parameters. And it does not order
finite matrices of different sizes: changing $N$ changes the compression, so
only matrices in the same subspace are comparable by this identity.

\subsection{The scalar table}

Each entry below is a certified enclosure on a \emph{fixed, frozen} function; no
reoptimization is involved, so the columns compare exact forms on the same
vector.

\begin{center}
\begin{tabular}{lll}
\toprule
 & $f_{200}$ & $f_{400}$ \\
\midrule
$R_{200}$ & $1.0310220302459896\times10^{-17}$ & $1.24275215\times10^{-17}$ \\
$R_{400}$ & $2.3520530467310079\times10^{-17}$ & $1.4330724047355850\times10^{-17}$ \\
$R_{800}$ & $3.01603\times10^{-17}$            & $1.6681\times10^{-17}$ \\
$Q$       & $3.6485917529040770\times10^{-17}$ & $1.9013829111454644\times10^{-17}$ \\
\bottomrule
\end{tabular}
\end{center}

The $R_{800}$ entries carry wider enclosures ($\pm 5.30\times10^{-23}$ and
$\pm 4.60\times10^{-22}$) because the quadrature remainder grows with the
cutoff; $R_{200}(f_{400})$ was evaluated as a direct scalar integral on the full
degree-$438$ function, with no truncation and no enlarged matrix. Each column is
strictly increasing in $T$, as Proposition~\ref{prop:cutoff} requires.

\subsection{What the comparison discards}

From~\eqref{eq:e}, with $e_T(f) = q(f) - r_T(f)$ by rigorous interval
subtraction:
\begin{center}
\begin{tabular}{lll}
\toprule
 & $f_{200}$ & $f_{400}$ \\
\midrule
$e_{200}$ & $2.617569722658087\times10^{-17}$ & $6.586308\times10^{-18}$ \\
$e_{400}$ & $1.296538706173069\times10^{-17}$ & $4.683105064098794\times10^{-18}$ \\
$e_{800}$ & $6.3256\times10^{-18}$            & $2.333\times10^{-18}$ \\
\midrule
$e_{400}/e_{200}$ & $0.49532156 \pm 4.69\times10^{-9}$ & $0.7110365 \pm 1.61\times10^{-8}$ \\
$e_{800}/e_{400}$ & $0.48788 \pm 9.07\times10^{-6}$    & $0.498 \pm 2.36\times10^{-4}$ \\
\bottomrule
\end{tabular}
\end{center}

\textbf{There is no decay law here.} On $f_{400}$ the successive ratios are
$0.711$ and then $0.498$, which is not a constant rate; and the two functions
disagree at the \emph{same} doubling $200 \to 400$, giving $0.495$ and $0.711$,
so the ratio is not a function of $T$ alone. That both functions give
approximately $0.49$ across $400 \to 800$ is two samples of one step, not a rate
in $T$, and comparing differently optimized vectors would establish less still.
No uniform operator estimate is inferred from a fixed-vector ratio, and none of
these ratios is used in any certified statement.

\subsection{What $q_f - r_f$ is, and is not}

For $f_{400}$, $q_f - r_f = 4.6831050640988\times10^{-18}$ is the
discarded-term identity~\eqref{eq:e} evaluated on that function, inside its
certified enclosure. It is \emph{not} $\lQ - \lR(400)$: both $\lR(400) \le r_f$
and $\lQ \le q_f$ are upper bounds, and two upper bounds do not make their
difference an equality between infima.

\subsection{Certified ceilings, and the stopping decision}

Because $f_{400}$ and $f_{200}$ are admissible, each $R_T$ value in the table is
a variational \emph{upper} bound for $\lR(T)$, and any certificate resting only
on $Q \ge R_{T,\beta}$ is capped by it. Relative to the lower bound in force:
\begin{center}
\begin{tabular}{llll}
\toprule
route & certified ceiling & reference lower bound & headroom \\
\midrule
$R_{200}$ & $\lR(200) \le r_{200}(f_{200}) = 1.0310220\times10^{-17}$ & $1\times10^{-17}$ & $3.102\%$ \\
$R_{400}$ & $\lR(400) \le r_{400}(f_{400}) = 1.4330724\times10^{-17}$ & $1.41\times10^{-17}$ & $1.636\%$ \\
$R_{800}$ & $\lR(800) \le r_{800}(f_{400}) = 1.6681\times10^{-17}$ & $1.41\times10^{-17}$ & $18.31\%$ \\
\bottomrule
\end{tabular}
\end{center}

So the $T = 400$ route is within about $1.6\%$ of exhausted, and the entire
$T = 800$ route is capped at a factor $1.1830$ over the present lower endpoint;
since a reoptimized $f_{800}$ can only sit below that ceiling, the realizable
gain is strictly less. Measured costs on the machine used:

\begin{center}
\begin{tabular}{ll}
\toprule
stage & measurement \\
\midrule
assembly, $T=200$, $N=220$, $800\times64$ nodes & $1338$ s, peak $205.5$ MiB \\
assembly, $T=400$, $N=440$, same node count    & $3532$ s, peak $393.4$ MiB \\
Rump factorizations, both sectors, $N=440$     & under $50$ s \\
compact $Q(f_{200})$                           & $389$ s \\
cross table ($R_{400}$, $Q(f_{400})$)          & $1544$ s \\
scalar leg $R_{200}(f_{400})$, $12800$ nodes   & $933$ s \\
scalar leg $R_{800}(f_{400})$, $51200$ nodes   & $8049$ s \\
scalar leg $R_{800}(f_{200})$, $51200$ nodes   & $3651$ s \\
\bottomrule
\end{tabular}
\end{center}

The assembly grew by a factor $2.64$ when $N$ doubled at fixed node count, and
the $T = 800$ scalar legs overran the node-count scaling by more than a factor
of two, which indicates that per-Bessel cost grows with the argument rather than
node count alone. Extrapolating --- as an estimate anchored on these
measurements, not a guarantee --- a full $T = 800$ matrix at $N \approx 880$ with
roughly $1600\times64$ nodes would cost on the order of ten hours of assembly.
Against a ceiling of $+18.31\%$, the recommendation recorded in the
accompanying repository is to \textbf{consolidate at $T = 400$}. Neither the
pilot nor any failed certificate here is an intrinsic barrier; this is a
cost--benefit judgement from measured data, and a separately scoped $T = 800$
attempt would be a reasonable thing to run with that ceiling stated in advance.

\section{Reproducibility and limitations}
\label{sec:repro}

\subsection{Artifacts and verification}

Every decisive quantity in this note is published as a midpoint--radius
enclosure in machine-readable form, together with the exact rational input
vectors and the code that produced them. Printed decimals in this note are
summaries; the enclosures are authoritative. Six independent verifiers
recompute the chain from the frozen artifacts and exit nonzero on any failure;
on the run accompanying this version all six exit $0$ with $119$ checks passing
in total ($34 + 19 + 20 + 16 + 18 + 12$), in $97$ seconds. A seventh verifier,
shipped with this note, re-derives every decimal quoted above from the
artifacts and additionally checks each quoted bound's \emph{rounding
direction} and its scope and parity label ($47$ further checks); it is what
caught the mis-quoted upper endpoint of Appendix~\ref{app:errors}, item~6.

The verifiers fail closed. Any NaN, infinite radius, indeterminate enclosure, or
unresolved sign produces an explicit failure or a refinement request; no skipped
cell, panel or pivot is interpreted as success. This matters because an earlier
version of the computation silently skipped cells where a library digamma call
returned NaN on wide balls and thereby produced a wrong constant that was
\emph{used} (Appendix~\ref{app:errors}, item~2).

\subsection{Environment}

Linux, Python $3.14.7$, \src{mpmath 1.3.0}, \src{python-flint 0.9.0} (FLINT/Arb).
Floating-point \src{mpmath} is used only to \emph{locate} candidate vectors,
which are then frozen as exact rationals and certified; it appears in no
pass/fail decision. Arb working precisions on sign-critical paths range from
$256$ to $2000$ bits and are stated per stage in the artifacts.

\subsection{Limitations}

\begin{enumerate}
\item \textbf{One window.} Everything is at $L = 4/5$. There is no propagation
  theorem, and nothing follows at any larger support. The repository's companion
  measurement at $L = 81/100$ indicates roughly $75\times$ the assembly work for
  a target that is certifiably at most $6.565\times10^{-18}$, which is why that
  window was not attempted here.
\item \textbf{An imported normalization.} The explicit-formula pairing in the
  convention of Section~\ref{sec:setup} is imported from the accompanying
  repository, not rederived. It was cross-checked only through the independent
  closed form for $Q[0,0]$ and the degree-$2$ convention checks of
  Section~\ref{sec:upper}.
\item \textbf{An unverified external matrix step.} Zhu describes supplementary
  source code and certificates in~\cite[\S17]{Zhu}; these were not obtained or
  executed in the present work, so the matrix step of~\cite[Thm.~1.1]{Zhu} is
  not reproduced here. The present certificate is independent of it.
\item \textbf{Shared arithmetic.} The compact time-domain evaluator shares
  library code with the matrix assembly, and is reported as a second
  representation, not an independent implementation.
\item \textbf{No arithmetic content.} The only analytic sharpening proved here
  is the monotone gate of Lemma~\ref{lem:gate}, which moves a constant by under
  one percent and does not touch the doubly exponential threshold that governs
  the cost in $L$.
\item \textbf{Scope.} No zero-free region, no verified range of zeros, no
  statement about the Riemann hypothesis.
\end{enumerate}

\subsection{Revision history}

\begin{itemize}
\item[\textbf{1.0.0}] (2026-11-10) First version.
\item[\textbf{1.0.1}] (2026-11-10) Citation and provenance revision only. Adds
  the two earlier technical notes whose implementation and parity conventions
  this computation inherits (\cite{Sensitivity} and~\cite{ParityNote}), the
  lineage paragraph in Section~\ref{sec:prior}, and the corresponding citations
  in Sections~\ref{sec:setup} and~\ref{sec:upper}. No numerical result, bound,
  artifact or limitation changed; the certified bracket and the parity
  statement are exactly those of version 1.0.0.
\item[\textbf{1.0.2}] (2026-11-10) Corrections after reading~\cite{Zhu} in full,
  which had previously been consulted only through its abstract and rendered
  HTML. Three attributions were too weak and are strengthened: \cite[Thm.~6.2]{Zhu}
  certifies the parity \emph{separation} at this very window, with a far stronger
  odd-sector constant, and establishes simplicity and evenness of the ground
  state (Sections~\ref{sec:prior} and~\ref{sec:parity});
  \cite[\S6.2]{Zhu} already evaluates its variational upper bounds by the
  time-domain route used in Section~\ref{sec:upper}; and the supplementary
  materials of~\cite[\S17]{Zhu} are described there, rather than absent.
  Section~\ref{sec:compared} is new and compares the two sets of bounds at this
  window. Again no numerical result, bound or artifact changed.
\end{itemize}

\appendix

\section{Error history}
\label{app:errors}

Errors found during this work are recorded because several of them produced
plausible wrong numbers rather than failures, and because the countermeasures
are part of what the artifacts are for. Items 1--5 are library- or
discipline-level defects; items 6--8 are corrections to claims. A longer ledger,
including retracted statements $R1$--$R13$, is kept with the artifacts.

\begin{enumerate}
\item \textbf{Pole-coordinate phase.} The form is assembled in the coordinates
  $y_n = \sigma_n x_n$ with $\sigma_n = (-1)^{\lfloor n/2\rfloor}$, so the pole
  vector requires the same phase. Omitting it is invisible at degree $0$ and
  wrong from degree $2$ onward. It produced a genuinely negative pivot
  ($-0.3949$) that would have been reported as indefiniteness. Caught by
  comparing $|F(t)|^2$ and $F(i/2)$ against direct integration; the degree-$2$
  checks of Section~\ref{sec:upper} now cover it.
\item \textbf{Digamma on wide balls.} The library digamma returns NaN on balls
  whose real part has radius $\gtrsim 0.1$, even strictly inside the right
  half-plane. A comparison loop silently skipped those cells and returned
  $M \approx 1831$ instead of $8.83$ --- and the wrong value was used. Replaced
  by the explicit recurrence $\psi(z) = \psi(z+K) - \sum_{k<K}(z+k)^{-1}$ with
  adaptive $K$, and by a covering-free monotonicity bound for $M$.
\item \textbf{Bessel functions on wide balls.} Interval evaluation of $j_n$ on
  the complex ellipses is useless here (width $\sim10^{32}$ at $|w|\sim80$,
  unchanged at $3000$ bits). Replaced by the uniform
  Cauchy--Schwarz bound~\eqref{eq:cs}.
\item \textbf{Interval factorization.} Naive interval $LDL^{\mathsf T}$ is
  unusable at these dimensions (radii amplify about $6.5\times$ per pivot,
  reaching $O(1)$ by step $47$). Replaced by the Rump preconditioned congruence.
\item \textbf{Precision ownership.} The Arb context precision is a process
  global. Routines set it on entry and did not restore it, so a callee silently
  changed its caller's context; this produced four separate wrong constants in
  this project, each plausible-looking rather than an error (a complement
  constant at $[\pm 7.4\times10^{132}]$, an even-sector $d$ at $[\pm
  2.8\times10^{183}]$, and two coefficient sets built at $53$ bits). The repair
  is an owned-scope wrapper that asserts restoration on normal return
  \emph{and} on exception, caches keyed by precision so a low-precision value
  can never be served to a high-precision caller, and an explicit
  \emph{unresolved} result for any wide or non-finite enclosure. The lesson, and
  the harder half, is that \emph{an input constructed outside the scope carries
  the ambient precision's radius, and entering a higher-precision scope
  afterwards does not recover it}: own your inputs as well as your scope.
\item \textbf{The upper endpoint, twice mis-quoted.} Two distinct errors, both
  in the safe direction numerically but both wrong as stated. (a) Internal
  drafts quoted the enclosure \emph{midpoint}
  $1.90138291114546440\times10^{-17}$ as the bracket's upper endpoint; a
  midpoint is not an upper bound for its own ball. (b) A separate record quoted
  $1.9013829111454646892\times10^{-17}$ as ``the exact upper endpoint''; that
  value came from re-parsing the serialized ball at the ambient $53$ bits,
  which rounds the midpoint and then inflates the endpoint by about
  $2.8\times10^{-33}$ --- the precision-ownership defect of item~5 recurring in
  the very routine that documents it. The strict endpoint is
  $U = 1.901382911145464406030000\times10^{-17}$, and
  $1.9014\times10^{-17}$ is the safe rounded form.
  See Remark~\ref{rem:endpoint}.
\item \textbf{A retracted sparsity target.} An earlier proposed next step --- a
  power-saving sparsity bound on the near-resonant set of the prime comb --- is
  false, and is retracted in full: by continuous Kronecker--Weyl equidistribution
  that set has positive density, confirmed numerically here to four significant
  figures across four decades of $T$. The replacement is the empty-exterior gate
  of Section~\ref{sec:lower}, which needs no such estimate.
\item \textbf{Overstated interpretations, corrected.} Four further claims made
  alongside earlier versions were withdrawn or corrected: that no bound on the
  measure of a bad set could control $\int_E|F|^2$ (false --- standard
  finite-measure and local-occupation estimates do exactly that); that a Schur
  \emph{threshold} was a lower bound (it is a break-even point, and the final
  test needs strict inequalities, hence~\eqref{eq:schur}); that lowering $\beta$
  would buy a larger $a$ (refuted by Proposition~\ref{prop:beta} --- the
  $\beta\|f\|_2^2$ term had been ignored); and an accusation that the library's
  adaptive integrator returned a silently wrong value (withdrawn --- it could not
  be reproduced; the wrong numbers came from a run that still carried the
  pole-phase error of item~1, at a different cutoff, with output from two runs
  conflated).
\end{enumerate}

\begin{thebibliography}{99}

\bibitem{Bombieri} E. Bombieri,
\emph{Problems of the Millennium: the Riemann Hypothesis},
Clay Mathematics Institute, 2000.

\bibitem{CC21} A. Connes and C. Consani,
\emph{Weil positivity and trace formula: the archimedean place},
Selecta Mathematica \textbf{27} (2021), no.~4, art.~77.

\bibitem{CCM} A. Connes, C. Consani and H. Moscovici,
\emph{Zeta zeros and prolate wave operators},
preprint, arXiv:2203.02261.

\bibitem{CvS} A. Connes and W. D. van Suijlekom,
\emph{Spectral truncations in noncommutative geometry and operator systems},
Communications in Mathematical Physics \textbf{383} (2021), 2021--2067.

\bibitem{Rump} S. M. Rump,
\emph{Verification methods: rigorous results using floating-point arithmetic},
Acta Numerica \textbf{19} (2010), 287--449.

\bibitem{Arb} F. Johansson,
\emph{Arb: efficient arbitrary-precision midpoint-radius interval arithmetic},
IEEE Transactions on Computers \textbf{66} (2017), no.~8, 1281--1292.

\bibitem{Slepian} D. Slepian and H. O. Pollak,
\emph{Prolate spheroidal wave functions, Fourier analysis and uncertainty I},
Bell System Technical Journal \textbf{40} (1961), 43--63.

\bibitem{Weil} A. Weil,
\emph{Sur les ``formules explicites'' de la th\'eorie des nombres premiers},
Comm. S\'em. Math. Univ. Lund (1952), 252--265.

\bibitem{Zhu} X. Zhu,
\emph{Weil positivity in compact windows: a finite reduction, certified
two-sided bounds, and a Landau--Widom decay law},
preprint, arXiv:2608.24827v2 (2026).

\bibitem{Sensitivity} E. De Jes\'us,
\emph{A $\delta^2$--$L^3$ Sensitivity Law for Finite-Window Weil Positivity,
with a Function-Field Control of the Connes--van Suijlekom Pipeline},
technical note, Zenodo, 2026.
\href{https://doi.org/10.5281/zenodo.23250932}{doi:10.5281/zenodo.23250932}

\bibitem{ParityNote} E. De Jes\'us,
\emph{Parity Blindness of Finite-Window Weil Detectors at a Real Zero:
the Even Sector Cannot See a Siegel Configuration and the Odd Sector Is
Maximally Sensitive to It},
technical note, version 2, Zenodo, 2026.
\href{https://doi.org/10.5281/zenodo.23624161}{doi:10.5281/zenodo.23624161}

\bibitem{Repo} E. De Jes\'us,
\emph{weil-window-control}, research repository, branch
\texttt{positive-weil-certificate-2026-11-09}.\\
\url{https://github.com/innerlightr-wq/weil-window-control}

\end{thebibliography}

\end{document}
